\section{Cross-Field Mechanism Map}
\label{sec:mechanisms}

Mechanisms attach to atomic conditions, not automatically to cumulative levels.
Table~\ref{tab:mapping} maps representative families by their direct
contribution and the assumption needed to compose that contribution. A row
marked G4 means conformance only to the statement or operator that the cited
construction actually verifies; it does not silently supply G1--G3.

\begin{table}[t]
\caption{Technique-to-goal map. A cumulative V$i$ claim requires every G1--G$i$
condition, possibly by composition. ``Direct'' denotes the condition to which
the technique principally contributes, not a universal theorem for the family.}
\label{tab:mapping}
\centering
\scriptsize
\setlength{\tabcolsep}{2pt}
\begin{tabular}{@{}L{2.0cm}L{2.25cm}L{1.6cm}L{5.4cm}@{}}
\toprule
Field / technique & Examples & Direct & Cumulative interpretation and caveat \\
\midrule
Tamper-evident logging & Merkle/history trees; hash chains
\cite{crosby2009} & G1 with an external witness & A producer-controlled root is
not G1. A tree binds leaves and a chain binds an ordered prefix; neither proves
coverage or availability of openings. \\

Record authentication & MACs, signatures, BLS aggregation
\cite{boneh2003aggregate} & G2 & Cumulative V2 also needs G1. It proves use of a
credential; device identity depends on enrollment and honest acquisition on the
capture profile. \\

Secure sensor aggregation & SIA
\cite{przydatek2003sia} & G1/G2 plus probabilistic,
operator-specific tests & Its sampling theorem and approximate aggregate must
be mapped to G3/G4 and $\delta$; approximate validity is not exact cumulative
V4. \\

Verifiable private aggregation & mPVAS; PPVDA
\cite{palazzo2024mpvas,zhou2024ppvda} & G2 and operator-specific G4 & G1 and G3
are not automatic. Public versus designated verification, collusion, supported
operator, and privacy assumptions remain explicit. \\

General verifiable computation & Groth16; proof-carrying data
\cite{groth2016,chiesa2010pcd} & G4 for the constrained statement & A circuit may
internalize G2/G3, but freshness and external anchoring still supply G1. Setup,
statement binding, and numerical semantics matter. \\

Attestation / proof of execution & RATS; APEX; VECODI
\cite{rfc9334,apex2020,vecodi2026} & execution identity and G4 under an isolated
path & Cumulative use needs appraisal, enrollment, input binding, coverage, and
a fresh external G1 event. Supports $\rho_{\mathrm{iso}}$, not physical truth. \\

Roots of trust / evaluation & Secure elements, TPMs, PUFs, evaluated hardware
& no G-condition alone & Supports a capture profile only for the evaluated key,
path, attacks, and duration; small report size is not a V-level. \\

Robust or statistical fusion & MATE and trust-aware fusion
\cite{mate2025} & physical-value plausibility / resilience &
Orthogonal to G1--G4 and useful precisely because computation-conformant V4
cannot establish physical truth. \\
\bottomrule
\end{tabular}
\end{table}

The map exposes several composition rules. Aggregate signatures compress tags,
not messages, identities, key bindings, or coverage. A homomorphic MAC can
authenticate a linear aggregate for a designated auditor while omitting the
manifest integrity and per-record evidence needed by a cumulative stack. A
succinct proof can instead constrain attribution and coverage inside its
statement, but its manifest commitment still needs a fresh G1 anchor.

As a worked mapping, SIA's authenticated query nonce and server-held Merkle
commitment contribute G1; sensor--server MACs on opened leaves contribute
probabilistic G2; duplicate, count, and input-representation tests contribute
probabilistic coverage under its ID-knowledge or alarm-channel assumptions;
and its interactive tests establish operator-specific approximation
\cite{przydatek2003sia}. It does not reach cumulative V3 here because G3 asks
for the exact policy-ordered identifier sequence, nor V4 because G4 asks for
exact $y=F(O_M)$ rather than an $\epsilon$-approximation. Its own
$(\epsilon,\delta)$ claim remains valuable and should be reported directly.

Sampling adds a protocol condition rather than a new semantic goal. The root
must be anchored before an unpredictable verifier challenge, and an offline
judge must compare the response with its own stored challenge or an
authenticated challenge transcript. Trusting a challenge copied from the
producer's proof permits root grinding. Public and designated verification,
and deterministic and probabilistic guarantees, must therefore be stated even
when they contribute to the same G-condition.

Proofs of execution occupy a different trust profile from cryptographic
arguments. Groth16 work follows the circuit \cite{groth2016}, whereas APEX binds
an output to execution under a hardware root \cite{apex2020}. In our modeled
attestation stack, an external authority first anchors the manifest root; the
isolated path then reports that root and $y$. The row measures only this
protocol transcript under $\rho_{\mathrm{iso}}$: it does not measure hardware,
energy, or physical resistance.
